\documentclass[a4paper,12pt]{article}

\usepackage[T2A]{fontenc}
\usepackage[utf8]{inputenc}
\usepackage[russian]{babel}
\usepackage{amsmath,amssymb,amsthm,mathtools}
\usepackage{helvet}
\renewcommand{\familydefault}{\sfdefault}
\usepackage{icomma}
\usepackage[a4paper,left=18mm,right=18mm,top=18mm,bottom=20mm]{geometry}
\usepackage{microtype}
\usepackage{tikz}
\usetikzlibrary{calc}
\usepackage{tkz-euclide}
\usepackage{pgfplots}
\pgfplotsset{compat=1.18}
\usepackage{tabularx,booktabs,longtable}
\usepackage{enumitem}
\usepackage{hyperref}
\usepackage{parskip}
\usepackage{fancyhdr}
\usepackage{setspace}
\usepackage{xcolor}

\definecolor{accent}{HTML}{008080}
\definecolor{darkaccent}{HTML}{004D4D}
\definecolor{textgray}{HTML}{333333}
\definecolor{softgray}{HTML}{6B7280}

\hypersetup{
  colorlinks=true,
  linkcolor=darkaccent,
  urlcolor=darkaccent
}

\color{textgray}
\onehalfspacing
\setlength{\parindent}{0pt}
\setlength{\headheight}{25pt}

\setlist[enumerate]{
  leftmargin=*,
  itemsep=0.35em,
  topsep=0.4em
}

\setlist[itemize]{
  leftmargin=*,
  itemsep=0.25em,
  topsep=0.35em
}

\pagestyle{fancy}
\fancyhf{}
\fancyhead[L]{\small\color{softgray}Математика \textbullet{} 7 класс}
\fancyhead[R]{\small\color{softgray}Кирилл Зиновьев}
\fancyfoot[C]{\small\color{softgray}\thepage}

\renewcommand{\headrulewidth}{0.3pt}
\renewcommand{\headrule}{
  \hbox to\headwidth{
    \color{accent}
    \leaders\hrule height \headrulewidth\hfill
  }
}

\newcommand{\sectionline}{
  \par
  \vspace{-0.35em}
  {\color{accent}\rule{\linewidth}{0.5pt}}
  \vspace{0.55em}
}

\newcommand{\skill}[1]{\textbf{\color{darkaccent}#1}}
\newcommand{\warn}[1]{\textbf{Важно:} #1}

\newcommand{\nodmath}{
  \mathop{\mbox{НОД}}\nolimits
}

\newcommand{\nokmath}{
  \mathop{\mbox{НОК}}\nolimits
}

\begin{document}

% =========================================================
% ТИТУЛЬНЫЙ ЛИСТ
% =========================================================

\begin{titlepage}
\thispagestyle{empty}
\centering

\vspace*{28mm}

{\Large\color{accent}\textbf{ЧЕК-ЛИСТ}}\\[9mm]

{\huge\bfseries
Повторение курса 6 класса
}\\[3mm]

{\Large\bfseries
Арифметика, НОД и НОК, дроби,\\
координатная прямая
}\\[16mm]

{\large
Кирилл Зиновьев, 7 класс
}\\[24mm]

{\large\color{darkaccent}
\textbf{Лёвин Артём Александрович}
}\\[2mm]

{\normalsize
эксклюзивно для Кирилла Зиновьева
}\\[12mm]

{\normalsize \today}

\vfill

\begin{tikzpicture}[remember picture,overlay]
  \draw[
    accent,
    line width=1.2pt
  ]
  ($(current page.south west)+(0,16mm)$)
  ..
  controls
  ($(current page.south west)+(55mm,29mm)$)
  and
  ($(current page.south east)+(-62mm,5mm)$)
  ..
  ($(current page.south east)+(0,17mm)$);
\end{tikzpicture}

\end{titlepage}

% =========================================================
% ОСНОВНОЙ ЧЕК-ЛИСТ
% =========================================================

\section*{1. Порядок действий в числовых выражениях}
\sectionline

\skill{Главный алгоритм.}
В выражении без дополнительных указаний действуйте в таком порядке:

\begin{enumerate}[label=\arabic*)]
  \item действия в скобках;
  \item степени, если они есть;
  \item умножение и деление слева направо;
  \item сложение и вычитание слева направо.
\end{enumerate}

Если в выражении есть несколько действий одного уровня,
нельзя произвольно менять их порядок.

Например,
\[
840:(36-22)+17\cdot 6
\]
начинают со скобок, затем выполняют деление и умножение,
после чего сложение.

\skill{Полезный приём для устного счёта.}
Перед вычислениями посмотрите, можно ли сгруппировать числа так,
чтобы получить целое число или удобную сумму.

Например,
\[
-3,7-5,2-(-1,8)-4,6
=
-3,7+(-5,2+1,8)-4,6.
\]

Здесь сначала удобно вычислить
\[
-5,2+1,8=-3,4.
\]

\warn{Сначала корректно раскройте скобки.}
В частности,
\[
-(-a)=+a.
\]

% =========================================================

\section*{2. Простые множители, НОД и НОК}
\sectionline

\skill{Разложение на простые множители}
--- представление натурального числа в виде произведения простых чисел.

Например,
\[
360=2^3\cdot 3^2\cdot 5,
\qquad
504=2^3\cdot 3^2\cdot 7.
\]

\skill{Признак делимости на 3.}
Число делится на $3$, если сумма его \emph{цифр}
делится на $3$.

Например,
\[
5+0+4=9,
\]
поэтому число $504$ делится на $3$.

\skill{НОД}
--- наибольший общий делитель нескольких чисел.

При разложении на простые множители для НОД берут только
общие простые множители в \emph{наименьших}
встречающихся степенях.

Для чисел $360$ и $504$:
\[
\nodmath(360,504)
=
2^3\cdot 3^2
=
72.
\]

Проверка:
\[
360:72=5,
\qquad
504:72=7.
\]

Следовательно, $72$ действительно делит оба числа.

\skill{НОК}
--- наименьшее общее кратное нескольких чисел.

Для НОК берут все простые множители,
которые встречаются хотя бы в одном из разложений,
в \emph{наибольших} встречающихся степенях.

Поэтому
\[
\nokmath(360,504)
=
2^3\cdot 3^2\cdot 5\cdot 7
=
2520.
\]

Проверка:
\[
2520:360=7,
\qquad
2520:504=5.
\]

Таким образом, $2520$ делится и на $360$, и на $504$.

\skill{Короткая смысловая опора:}
\[
\boxed{
\text{НОД делит данные числа, а НОК делится на данные числа.}
}
\]

Для трёх чисел алгоритм остаётся тем же.

Например,
\[
8=2^3,
\qquad
12=2^2\cdot 3,
\qquad
15=3\cdot 5.
\]

Тогда
\[
\nokmath(8,12,15)
=
2^3\cdot 3\cdot 5
=
120,
\]
а
\[
\nodmath(8,12,15)=1.
\]

Общего простого множителя у всех трёх чисел нет,
поэтому их единственный общий положительный делитель,
подходящий для НОД, равен $1$.

\warn{Фраза «в НОК берём уникальные множители» недостаточно точна.}
Надёжное правило: для каждого простого множителя
берите максимальную степень, встречающуюся в разложениях.

% =========================================================

\section*{3. Сравнение обыкновенных и десятичных дробей}
\sectionline

Чтобы сравнить дроби разных видов, удобно привести их
к общему формату.

\skill{Шаг 1.}
При необходимости переведите десятичную дробь
в обыкновенную:

\[
0,6
=
\frac{6}{10}
=
\frac35.
\]

\skill{Шаг 2.}
Найдите общий знаменатель.

Наименьший общий знаменатель нескольких дробей ---
это НОК их знаменателей.

Рассмотрим числа
\[
\frac7{12},
\qquad
\frac35,
\qquad
\frac{11}{18},
\qquad
\frac58.
\]

Разложим знаменатели:
\[
12=2^2\cdot 3,
\qquad
5=5,
\]
\[
18=2\cdot 3^2,
\qquad
8=2^3.
\]

Следовательно,
\[
\nokmath(12,5,18,8)
=
2^3\cdot 3^2\cdot 5
=
360.
\]

\skill{Шаг 3.}
Приведите каждую дробь к знаменателю $360$:

\[
\frac7{12}
=
\frac{7\cdot 30}{12\cdot 30}
=
\frac{210}{360},
\]

\[
\frac35
=
\frac{3\cdot 72}{5\cdot 72}
=
\frac{216}{360},
\]

\[
\frac{11}{18}
=
\frac{11\cdot 20}{18\cdot 20}
=
\frac{220}{360},
\]

\[
\frac58
=
\frac{5\cdot 45}{8\cdot 45}
=
\frac{225}{360}.
\]

Теперь достаточно сравнить числители:
\[
210<216<220<225.
\]

Получаем
\[
\frac7{12}
<
0,6
<
\frac{11}{18}
<
\frac58.
\]

\warn{Дополнительный множитель применяется одновременно
к числителю и знаменателю.}

Например, чтобы перейти от знаменателя $12$ к $360$,
нужно вычислить
\[
360:12=30.
\]

Следовательно,
\[
\frac7{12}
=
\frac{7\cdot30}{12\cdot30}.
\]

% =========================================================

\section*{4. Рациональные числа и удобные вычисления}
\sectionline

\skill{Раскрытие скобок.}
Если перед скобками стоит знак минус,
знак каждого слагаемого внутри скобок меняется
на противоположный:

\[
a-(b-c)=a-b+c.
\]

В частности,
\[
-(-1,8)=+1,8.
\]

\skill{Сложение чисел разных знаков.}
Из большего модуля вычитают меньший
и сохраняют знак числа с большим модулем.

Например,
\[
-5,2+1,8=-3,4.
\]

Поскольку
\[
5,2>1,8,
\]
результат будет отрицательным.

\skill{Удобная группировка.}
Коммутативность и сочетательность сложения позволяют
переставлять и группировать слагаемые:

\[
a+b=b+a,
\]

\[
(a+b)+c=a+(b+c).
\]

Это особенно полезно, если некоторые слагаемые дают:

\begin{itemize}
  \item ноль;
  \item целое число;
  \item число с нулём в конце;
  \item другую простую для вычисления сумму.
\end{itemize}

% =========================================================

\section*{5. Координатная прямая, модуль, расстояние и середина}
\sectionline

Пусть на координатной прямой заданы точки
\[
A(-4,5),
\qquad
B(2,5).
\]

\begin{center}
\begin{tikzpicture}[x=1.15cm,y=1cm]

  \draw[->,line width=0.9pt]
    (-5.7,0)--(4.3,0)
    node[right] {$x$};

  \foreach \x in {-5,-4,-3,-2,-1,0,1,2,3,4}{
    \draw (\x,0.10)--(\x,-0.10);
    \node[below=4pt] at (\x,0) {\small $\x$};
  }

  \fill[accent] (-4.5,0) circle (2.2pt);
  \fill[accent] (2.5,0) circle (2.2pt);
  \fill[darkaccent] (-1,0) circle (2.2pt);

  \node[above=5pt] at (-4.5,0) {$A(-4,5)$};
  \node[above=5pt] at (2.5,0) {$B(2,5)$};
  \node[above=5pt] at (-1,0) {$C(-1)$};

\end{tikzpicture}
\end{center}

\skill{Выбор масштаба.}
Если координаты содержат половины единицы,
удобно выбрать такой масштаб, чтобы $0,5$
соответствовало целому числу клеток.

Например, можно принять две клетки за одну единицу.

\skill{Модуль числа}
--- расстояние от точки с данной координатой до нуля.

Поэтому
\[
|-4,5|=4,5.
\]

Модуль любого числа неотрицателен:
\[
|a|\ge 0.
\]

Для положительного числа
\[
|2,5|=2,5.
\]

Для отрицательного числа
\[
|-2,5|=2,5.
\]

\skill{Расстояние между двумя точками}
на координатной прямой рассчитывается по формуле

\[
AB=|x_B-x_A|.
\]

Для точек $A(-4,5)$ и $B(2,5)$:
\[
AB
=
|2,5-(-4,5)|.
\]

Раскрываем скобки:
\[
AB
=
|2,5+4,5|
=
|7|
=
7.
\]

Если заранее известно, какая координата больше,
можно использовать правило

\[
\text{расстояние}
=
\text{большая координата}
-
\text{меньшая координата}.
\]

То есть
\[
AB
=
2,5-(-4,5)
=
7.
\]

\skill{Координата середины отрезка}
вычисляется по формуле

\[
x_C
=
\frac{x_A+x_B}{2}.
\]

Для данных точек
\[
x_C
=
\frac{-4,5+2,5}{2}.
\]

Получаем
\[
x_C
=
\frac{-2}{2}
=
-1.
\]

Следовательно, середина отрезка $AB$ ---
точка
\[
C(-1).
\]

\skill{Проверка середины.}
Середина должна быть равноудалена от концов отрезка:

\[
AC
=
|-1-(-4,5)|
=
3,5,
\]

\[
CB
=
|2,5-(-1)|
=
3,5.
\]

Следовательно,
\[
AC=CB.
\]

\warn{При подстановке отрицательной координаты сохраняйте скобки.}

Записи
\[
2,5-(-4,5)
\]
и
\[
2,5-4,5
\]
имеют разные значения.

% =========================================================

\section*{6. Итоговый чек-лист перед входной диагностикой}
\sectionline

Проверьте, умеете ли Вы без подсказки:

\begin{enumerate}[label=\arabic*)]

  \item определять правильный порядок действий
  в числовом выражении;

  \item раскрывать скобки и корректно работать
  со знаками;

  \item раскладывать натуральные числа
  на простые множители;

  \item применять признаки делимости,
  особенно на $2$, $3$, $5$, $9$, $10$;

  \item находить НОД через минимальные степени
  общих простых множителей;

  \item находить НОК через максимальные степени
  всех необходимых простых множителей;

  \item переводить конечную десятичную дробь
  в обыкновенную;

  \item сокращать обыкновенные дроби;

  \item находить общий знаменатель через НОК
  знаменателей;

  \item приводить дроби к общему знаменателю;

  \item сравнивать дроби после приведения
  к общему знаменателю;

  \item выполнять действия с положительными
  и отрицательными числами;

  \item использовать удобную группировку
  слагаемых при вычислениях;

  \item находить модуль числа;

  \item вычислять расстояние между точками
  на координатной прямой;

  \item находить координату середины отрезка
  по формуле среднего арифметического координат.

\end{enumerate}

% =========================================================
% ТРЕНИРОВОЧНЫЙ БЛОК
% =========================================================

\clearpage

\section*{Тренировочный блок}
\sectionline

\textbf{Формат: Путь А.}

Выполните 10 заданий без подсказок.
Задания расположены по нарастающей сложности.

\begin{enumerate}[label=\textbf{\arabic*.},leftmargin=*]

  % ---------- Средний уровень ----------

  \item
  Вычислите:
  \[
  960:(52-28)+14\cdot 7.
  \]

  \item
  Разложите числа $420$ и $630$
  на простые множители и найдите
  \[
  \nodmath(420,630),
  \qquad
  \nokmath(420,630).
  \]

  \item
  Точки имеют координаты
  \[
  A(-3,5),
  \qquad
  B(4,5).
  \]

  Найдите:

  \begin{enumerate}[label=\arabic*)]
    \item $|-3,5|$;
    \item расстояние $AB$;
    \item координату середины отрезка $AB$.
  \end{enumerate}

  % ---------- Выше среднего ----------

  \item
  Найдите
  \[
  \nokmath(18,24,30)
  \]
  и
  \[
  \nodmath(18,24,30).
  \]

  \item
  Расположите в порядке возрастания числа
  \[
  \frac{5}{12},
  \qquad
  0,45,
  \qquad
  \frac{7}{15},
  \qquad
  \frac12.
  \]

  \item
  Вычислите наиболее удобным способом:
  \[
  -6,8-3,7-(-1,8)+0,7.
  \]

  \item
  Найдите наименьший общий знаменатель дробей
  \[
  \frac{5}{18},
  \qquad
  \frac{7}{24},
  \qquad
  \frac{11}{30},
  \qquad
  \frac{13}{40},
  \]
  а затем расположите эти дроби
  в порядке возрастания.

  \item
  Точки
  \[
  M(-7,2)
  \quad\text{и}\quad
  N(1,6)
  \]
  лежат на координатной прямой.

  Найдите расстояние $MN$
  и координату середины отрезка $MN$.

  Затем проверьте, что найденная середина
  равноудалена от $M$ и $N$.

  \item
  Найдите натуральное число $n$, если
  \[
  \nodmath(72,n)=18,
  \qquad
  \nokmath(72,n)=360.
  \]

  Укажите все возможные значения $n$.

  % ---------- Сложное ----------

  \item
  Четыре числа
  \[
  \frac{11}{24},
  \qquad
  0,47,
  \qquad
  \frac{7}{15},
  \qquad
  \frac{23}{50}
  \]
  нужно расположить в порядке возрастания.

  Выполните сравнение точно, без округления.

  Затем найдите расстояние между наименьшим
  и наибольшим из этих чисел,
  если рассматривать их как координаты точек
  на числовой прямой.

\end{enumerate}

% =========================================================
% ОТВЕТЫ
% =========================================================

\clearpage

\section*{Ответы к тренировочному блоку}
\sectionline

\renewcommand{\arraystretch}{1.35}

\begin{table}[ht]
\centering
\caption{Краткие ответы}

\begin{tabularx}{\linewidth}{@{}>{\bfseries}p{12mm}X@{}}

\toprule
№ & Ответ \\
\midrule

1 &
$138$.
\\

2 &
\[
420=2^2\cdot3\cdot5\cdot7,
\qquad
630=2\cdot3^2\cdot5\cdot7.
\]
\[
\nodmath(420,630)=210,
\qquad
\nokmath(420,630)=1260.
\]
\\

3 &
1) $3,5$;
2) $8$;
3) $0,5$.
\\

4 &
\[
\nokmath(18,24,30)=360,
\qquad
\nodmath(18,24,30)=6.
\]
\\

5 &
\[
\frac{5}{12}
<
0,45
<
\frac{7}{15}
<
\frac12.
\]
\\

6 &
$-8$.
\\

7 &
Наименьший общий знаменатель:
\[
360.
\]
Порядок:
\[
\frac{5}{18}
<
\frac{7}{24}
<
\frac{13}{40}
<
\frac{11}{30}.
\]
\\

8 &
\[
MN=8,8.
\]
Координата середины:
\[
-2,8.
\]
Расстояние от середины до каждого конца:
\[
4,4.
\]
\\

9 &
\[
n=90.
\]
\\

10 &
\[
\frac{11}{24}
<
\frac{23}{50}
<
\frac{7}{15}
<
0,47.
\]

Расстояние между крайними числами:
\[
0,47-\frac{11}{24}
=
\frac{7}{600}.
\]
\\

\bottomrule

\end{tabularx}
\end{table}

\end{document}